\begin{proof}
Ya se sabe por (5.11.6) que \emph{a)} y \emph{b)} son equivalentes, al igual que
\emph{c)} y \emph{d)}, y que \emph{a)} implica \emph{d)}. En el caso presente, la equivalencia
de \emph{a)} y \emph{d)} resulta de aplicar (7.2.3) a un anillo cociente íntegro
$B$ de $A$: la condición \emph{d)} es una formulación equivalente de la
condición (7.2.3, \emph{b)}) en virtud de (5.9.3.1). La equivalencia de \emph{e)}
y \emph{f)} es consecuencia de aplicar (7.2.2) al propio
$\mathscr O_Y$-módulo coherente $\mathscr O_Y$. Ya se ha observado
(5.11.7, (ii)) que, si $A$ satisface \emph{a)}, también la satisface toda
$A$-álgebra finita y todo anillo de fracciones de $A$. Por tanto,
la condición \emph{a)} para $A$ implica (con las notaciones de \emph{e)}) que
$C$ satisface \emph{a)} y, en consecuencia, también \emph{c)}. Como $C$ es
íntegro y $\dim(C)\geqslant
2$, se puede aplicar a $C$ la condición
\emph{c)} tomando como $T$ el punto cerrado de $Y=\operatorname{Spec}(C)$;
esto prueba que \emph{a)} implica \emph{e)}. Queda por demostrar que \emph{e)}
implica \emph{a)}; evidentemente, basta tratar el caso en que $A$ es íntegro y $B=A$.
Entonces hay que mostrar que la condición \emph{e)} implica la condición
(7.2.3, \emph{c)}). Con las notaciones de (7.2.3, \emph{c)}), consideremos un punto
$y'\in f^{-1}(T)\cap\overline{\{x'\}}$ y pongamos $y=f(y')$
y $C=\mathscr O_{X,y}$. Como $y\in T$, por hipótesis se tiene
$\dim(C)\geqslant2$ y, por consiguiente (con las notaciones de \emph{e)}),
$j_*(\mathscr O_U)$ es un $\mathscr O_Y$-módulo coherente. Pongamos
$Y'=X'\times_XY$. Como $f:X'\to X$ es plano, también lo es
$g=f_{(Y)}:Y'\to Y$. Además, el espacio subyacente a $Y'$ se
identifica con $f^{-1}(Y)$ (\textbf I, 3.6.5). Como $A$ es íntegro y $f$
es plano, $\operatorname{Ass}(\mathscr O_{X'})$ está contenido en
la fibra del punto genérico de $X$ (3.3.2); dicho punto pertenece
a $Y$, de donde $x'\in Y'$. Sea $U'=g^{-1}(U)$ y sea $j'$ la
inclusión canónica $U'\to Y'$. Puesto que
$j_*(\mathscr O_U)$ es un $\mathscr O_Y$-módulo coherente y $g$
es plano, de (5.9.4) resulta que $j'_*(\mathscr O_{U'})$ es un
$\mathscr O_{Y'}$-módulo coherente. Ahora bien,
$x'\in\operatorname{Ass}(\mathscr O_{U'})$, por lo que de
(5.10.10) se concluye que, en $Y'$,

\oldpage[IV-2]{190}

$\operatorname{codim}(\overline{\{y'\}}\cap\overline{\{x'\}},
\overline{\{x'\}})\geqslant2$.

Como $y'$ es arbitrario en
$f^{-1}(T)\cap\overline{\{x'\}}$, se obtiene
$\operatorname{codim}(f^{-1}(T)\cap\overline{\{x'\}},
\overline{\{x'\}})\geqslant2$ en $X'$. Queda demostrado.
\end{proof}
